\section{Related Literature}
\label{sec:related}

\subsection{Factor replication: HXZ, C\&Z, and JKP}

This paper joins a decade-long methodological debate about how to replicate
the published cross-sectional return-predictability literature. \citet{HouXueZhang2020}
rerun 452 anomalies under one uniform protocol. Over the original-study
samples, 65.3\% fail to clear $|t|\geq 1.96$ with NYSE breakpoints and value
weighting; the failure rate falls to 43.1\% under an equal-weighted protocol
with NYSE-Amex-NASDAQ breakpoints, a fact HXZ themselves use to
argue that weighting and breakpoint conventions, not the underlying
economics, drive much of the disagreement. Their contribution is an
aggregate failure rate under a single, externally imposed protocol; it does
not decompose any individual factor's gap into a signal component and a
configuration component, nor does it compare against the paper's own
originally reported number field by field.

\citet{ChenZimmermann2022} take the opposite approach: for each of 319
characteristics, they implement the paper's own stated method as closely as
possible, and report that all but a handful replicate. Where their result
disagrees with HXZ's, they make a manual, case-by-case judgment of whether
the anomaly itself is invalid or whether the standardized protocol departs
from what the paper actually specified. This is a substantive contribution,
but it stops short of an actual counterfactual: C\&Z never run their own
signal through HXZ's protocol, or vice versa, to isolate how much of the
disagreement is attributable to the signal formula versus the portfolio
configuration.

\citet{JensenKellyPedersen2023} instead replace the entire protocol with
a uniform Bayesian framework and report a sequence of six failure rates
(35\% $\to$ 55.6\% $\to$ 61.3\% $\to$ 82.4\% $\to$ 75.6\% $\to$ 82.4\%) as
successive methodological choices are changed one at a time -- comparable
tercile/decile conventions, exclusion of originally-insignificant anomalies,
CAPM-alpha rather than raw returns, and multiple-testing corrections. This
is the closest existing work to a decomposition, but it decomposes the
effect of \emph{replacing} the protocol, not the disagreement between two
independently-arrived-at implementations of the \emph{same} protocol space.

Table~\ref{tab:related-work} summarizes this positioning. What none of the
three provides is a genuine counterfactual execution -- taking one signal
and running it under another team's configuration, holding everything else
fixed -- that would cleanly separate "the signal is measured differently"
from "the portfolio is constructed differently." The $\Acz$/$\Ahxz$ hybrid
tracks in this paper (\S\ref{sec:design}) are built to do exactly that.

% Positioning table vs. HXZ (2020) / C&Z (2021-2022) / JKP (2023).
% See docs/paper-outline.md Ch.2 for the PDF-verified numbers (2026-08-18).
\begin{table}[htbp]
  \centering
  \caption{Positioning relative to HXZ, C\&Z, and JKP.}
  \label{tab:related-work}
  \begin{tabular}{p{2.2cm}p{5.5cm}p{5.5cm}}
    \toprule
    Source & What it does & What is missing, and what this paper adds \\
    \midrule
    HXZ (2020) & 452 anomalies, uniform protocol (NYSE breakpoints + VW); 65\% fail at $|t|\ge 1.96$; 43.1\% under EW + NYSE-Amex-NASDAQ breakpoints & Only an aggregate failure rate; no per-factor field-level gap \\
    C\&Z (2021/2022) & Manual case-by-case judgment of factor-invalid vs.\ method-deviates-from-paper; $\sim$100\% replication on 319 characteristics & Manual classification; never actually runs the C\&Z signal through the HXZ configuration as a counterfactual \\
    JKP (2023) & Uniform Bayesian protocol rerun; Figure 1 six bars (35\%$\to$55.6\%$\to$61.3\%$\to$82.4\%$\to$75.6\%$\to$82.4\%) attributed to protocol choices & Flattens disagreement (replaces the protocol) rather than decomposing it \\
    This paper & $A_{cz}$/$A_{hxz}$ hybrid tracks: genuine counterfactual execution & Cleanly splits each factor's total gap into signal- and configuration-difference components \\
    \bottomrule
  \end{tabular}
\end{table}


A separate, older literature frames the same tension without proposing a
decomposition: \citet{HarveyLiuZhu2016} and \citet{McLeanPontiff2016}
both treat the multiplicity of possible implementation and testing choices
as a source of noise to be corrected for (via higher significance
thresholds or out-of-sample decay, respectively) rather than as an object of
study in its own right. Complementary work addresses the high-dimensional
factor zoo by testing whether a proposed factor adds information beyond the
large set already documented \citep{FengGiglioXiu2020}. This paper's departure
is to treat exactly that
multiplicity -- the space of choices a paper's text leaves open -- as the
primary object of measurement.

\subsection{Researcher degrees of freedom and multiverse analysis}

The idea that a single dataset admits many defensible analytical choices,
and that the choice matters, is well established outside finance.
Specification-curve analysis \citep{Simonsohn2020} and
multiverse analysis \citep{Steegen2016} both enumerate the full space of
reasonable choices a researcher \emph{could} make and report how the result
varies across that space. Many-analyst evidence likewise shows that defensible
choices by independent teams can generate materially different results from
the same dataset and research question \citep{SilberzahnEtAl2018}. The finance
literature's own version of this concern is the p-hacking / multiple-testing
critique of anomaly discovery.

This paper's disagreement band is narrower by construction, and that is a
deliberate design choice rather than a limitation to apologize for: rather
than enumerate every choice a researcher could imaginably make, it is
calibrated to the choices two real, independent, well-resourced
implementers (the agent and Chen-Zimmermann) actually disagreed on. The
resulting band is an empirically observed interval in return outcomes, not an
a priori enumeration of hypothetical alternatives. Conditional on both
endpoints being admissible readings of the paper, its width is a lower bound
on the range of returns generated by the paper's implementation space
(\S\ref{sec:framework}) -- it is anchored to what real disagreement looks
like, not to what disagreement could in principle look like.

\subsection{LLM agents for research automation}

A growing literature uses large language models to automate parts of the
empirical research pipeline, including literature work, data analysis, coding,
and mathematical derivations \citep{Korinek2023}. The agent literature adds
interleaved reasoning and action \citep{YaoEtAl2023ReAct}, external-tool use
\citep{SchickEtAl2023Toolformer}, and interfaces through which a model edits
code and executes tests \citep{YangEtAl2024SWEAgent}. A parallel concern is
that an LLM's output can silently encode choices the model was never
authorized, or expected, to make. This paper's design responds to that concern
directly rather than by assertion: the LLM is only ever permitted to author
the \texttt{compute\_signal} function (\S\ref{sec:pipeline}); it never sets an
empirical parameter -- weighting, breakpoint convention, sample window, or
estimator -- all of which are drawn from a fixed, human-auditable configuration
menu. Every quantity reported in \S\ref{sec:results} is therefore reproducible
with the LLM switched off, once a plugin has been frozen; the LLM's role is
confined to producing an artifact whose downstream consequences are then
measured by a deterministic engine, not to producing the empirical findings
themselves. Persisting the specification, code, data provenance, and outputs
also follows established computational-reproducibility principles
\citep{StoddenEtAl2016}.

% estimator) -- those come from a fixed, human-auditable config menu.
